% Part: computability
% Chapter: computability-theory
% Section: russells-paradox

\documentclass[../../../include/open-logic-section]{subfiles}

\begin{document}

\olfileid[fr]{cmp}{thy}{rus}
\olsection{Comparaison avec le paradoxe de Russell}

Il est instructif de comparer les arguments
de cette section au paradoxe de Russell :
\begin{enumerate}
\item Le paradoxe de Russell : posons
  $S = \Setabs{x}{x \notin x}$. Alors
  $S\in S$ si et seulement si $S\notin S$,
  ce qui est contradictoire.\footnote{Correction éditoriale : la source écrit $X\notin S$ avec un $X$ non défini ; la condition appliquée à $S$ donne bien $S\notin S$.}

  \emph{Conclusion :} un tel ensemble~$S$
  n'existe pas. Supposer l'existence d'un
  « ensemble de tous les ensembles » est
  incompatible avec les axiomes usuels de
  la théorie des ensembles, par exemple ceux de ZF.

\item Une variante du paradoxe de Russell :
  soit $F$ la « fonction » définie de
  l'ensemble de toutes les fonctions vers
  $\{0,1\}$ par
  \[
  F(f) =
  \begin{cases}
    1 & \text{si $f$ appartient au domaine de $f$ et $f(f)=0$}\\
    0 & \text{sinon}
  \end{cases}
  \]
  Un raisonnement semblable montre que
  $F(F)=0$ si et seulement si $F(F)=1$ :
  contradiction.

  \emph{Conclusion :} $F$ n'est pas une
  fonction au sens ensembliste. L'« ensemble de toutes les
  fonctions » est trop grand pour être
  le domaine d'une fonction.

\item L'argument diagonal : soient
  $f_0,f_1,\dots$ l'énumération des fonctions
  partielles calculables, et
  $G\colon\Nat\to\{0,1\}$ la fonction définie par
  \[
  G(x) =
  \begin{cases}
    1 & \text{si $f_x(x)\downarrow = 0$}\\
    0 & \text{sinon}
  \end{cases}
  \]
  Si $G$ était calculable, elle coïnciderait
  avec une fonction $f_k$ de la liste.
  Mais alors $G(k)=1$ si et seulement si
  $G(k)=0$ : contradiction.

  \emph{Conclusion :} $G$ n'est pas calculable.
  Selon les axiomes de la théorie des ensembles,
  $G$ reste néanmoins une fonction : il n'y
  a pas ici de paradoxe, mais une distinction
  à clarifier.
\end{enumerate}

Les expressions « fonction partielle »,
« fonction calculable » et « fonction
partielle calculable » peuvent prêter à
confusion. L'ensemble de toutes les
fonctions partielles de $\Nat$ dans
$\Nat$ forme une vaste collection.
Certaines sont totales, certaines
calculables, certaines possèdent les
deux propriétés, et d'autres aucune.
Rappelons que, sans autre précision,
« fonction » désigne ici une fonction
totale. Nous distinguons donc :
\begin{enumerate}
\item les fonctions calculables;
\item les fonctions partielles calculables
  qui ne sont pas totales;
\item les fonctions qui ne sont pas calculables;
\item les fonctions partielles qui ne sont
  ni totales ni calculables.
\end{enumerate}
Pour visualiser ces cas, on peut dessiner
un grand carré représentant les fonctions
partielles de $\Nat$ dans $\Nat$, puis
y délimiter deux régions qui se recoupent :
celle des fonctions totales et celle des
fonctions partielles calculables. Chercher
un exemple dans chacune des régions du
diagramme est un bon exercice.

\end{document}
